\section{2.3 밴딧 vs 완전한 MDP: 다음 장으로 가는 다리}
\begin{frame}{밴딧에는 ``상태''가 없다}
  \begin{itemize}
    \item 팔을 당겨도 ``다음 상태''로 넘어가지 \textbf{않음}.
    \item 매번 같은 $k$개의 팔 중에서 고르는, 시간이 흘러도
      상황이 \textbf{전혀 안 바뀌는} 세계.
  \end{itemize}
  \vskip0.5em
  Week 3의 \kb{MDP(Markov Decision Process)}는 여기에 결정적인
  것 하나를 추가:
  \vskip0.3em
  \begin{center}
    \kq{``지금 한 행동이 다음에 어떤 \textbf{상태}를 보게 될지에
    영향을 준다.''}
  \end{center}
  \vskip0.3em
  슬롯머신 vs 체스: 체스에서는 지금 둔 수가 \textbf{판의 배치(상태)}
  를 바꾸고, 그 배치가 다음에 무엇을 할 수 있는지를 결정.
  \vskip0.3em
  밴딧에서 배운 탐험-활용 딜레마는 그대로 남지만,
  ``지금 선택이 \textbf{미래의 선택지 자체}를 바꾼다''는
  완전히 새로운 문제가 겹쳐짐.
\end{frame}
